\chapter{Estado del Arte}
\label{cap:estado}

Este capítulo resume la base técnica sobre la que se construye el sistema. Se revisan los LLMs y el motivo por el que, por sí solos, no son suficientes en entornos corporativos; se explica RAG como arquitectura práctica para responder con evidencia; y se describen los componentes clave para ejecutar todo en local: embeddings, bases vectoriales, cuantización, LoRA y MCP como estándar emergente de integración de herramientas.

% ============================================================
\section{Grandes Modelos de Lenguaje (LLMs)}
% ============================================================
Los grandes modelos de lenguaje (LLMs) son modelos neuronales entrenados para predecir el siguiente \textit{token} en una secuencia. El salto de rendimiento moderno llega con la arquitectura \textit{Transformer} \cite{vaswani2017attention}, donde la atención permite capturar relaciones de contexto de forma mucho más efectiva que arquitecturas previas.

\subsection{Propietarios vs. pesos abiertos}
En la práctica, hay dos enfoques:

\begin{itemize}
    \item \textbf{Modelos propietarios vía API (SaaS):} ofrecen rendimiento alto y producto maduro, pero implican dependencia de un proveedor externo y un flujo de datos fuera del perímetro. Además, normalmente no se publica información completa sobre el modelo (tamaño exacto, detalles de entrenamiento o parámetros).
    \item \textbf{Modelos de pesos abiertos (\textit{open-weights}):} permiten ejecución local u \textit{on-premise}, lo que habilita control de datos, trazabilidad y operación sin conectividad externa. Ejemplos representativos son Llama 3.1 \cite{meta2024llama31,meta2024llama31hf} y Qwen2.5 \cite{qwen2024qwen25blog}.
\end{itemize}

En este proyecto se prioriza el segundo enfoque porque el requisito principal es operar en local. Es importante distinguir ``pesos abiertos'' de ``código abierto'': la disponibilidad de pesos no implica licencias OSI y puede incluir condiciones de uso específicas según el proveedor \cite{meta2024llama31hf,google2024gemmaterms,qwen2024qwen25license}.

Se seleccionan Llama 3.1 8B y Qwen2.5 por ser viables en despliegue local y por su ecosistema \cite{meta2024llama31,qwen2024qwen25blog}.

% ============================================================
\section{Generación Aumentada por Recuperación (RAG)}
% ============================================================
Un LLM no consulta bases de datos en tiempo real durante la generación: produce texto a partir de patrones aprendidos. En un entorno corporativo eso es un problema porque la respuesta debe estar respaldada por documentación actual y verificable. Para resolverlo se utiliza RAG, que introduce una etapa de recuperación antes de la generación \cite{lewis2020rag}.

Un pipeline RAG típico tiene dos fases:

\begin{enumerate}
    \item \textbf{Retrieval:} la consulta se transforma en un embedding y se recuperan fragmentos relevantes desde un índice vectorial mediante búsqueda de vecinos próximos (ANN).
    \item \textbf{Generation:} el LLM recibe la pregunta junto con el contexto recuperado. Esto reduce respuestas no fundamentadas y permite devolver trazabilidad hacia los documentos usados.
\end{enumerate}

\subsection{Mejoras prácticas: híbrida y reranking}
En producción, RAG casi nunca se deja ``simple'':

\begin{itemize}
    \item \textbf{Búsqueda híbrida:} se combina recuperación densa (embeddings) con recuperación dispersa (BM25) para no perder coincidencias exactas (IDs, nombres de campos, artículos legales, acrónimos) \cite{robertson2009bm25,weaviateHybridDocs,qdrantHybridTutorial}. Esta es la variante implementada en el proyecto.
    \item \textbf{Re-ranking:} se reordenan los fragmentos candidatos con un modelo o heurística adicional para seleccionar los más relevantes antes de construir el contexto.
    \item \textbf{Memoria conversacional acotada:} se añade contexto del historial cuando aporta valor, evitando que crezca el \textit{prompt} sin control.
\end{itemize}

% ============================================================
\section{Modelos de Embeddings}
% ============================================================
Los embeddings convierten texto en vectores numéricos donde la distancia representa similitud semántica. En RAG, los embeddings son la base del retrieval. Una métrica común es la similitud del coseno:

\begin{equation}
\text{sim}(\mathbf{a}, \mathbf{b}) = \frac{\mathbf{a} \cdot \mathbf{b}}{||\mathbf{a}|| \cdot ||\mathbf{b}||}
\end{equation}

Para embeddings en local, una opción práctica y extendida es SentenceTransformers (Sentence-BERT) \cite{reimers2019sentence,sbertDocs}. Frente a alternativas comerciales (p.ej. OpenAI embeddings) \cite{openaiEmbeddingsGuide}, estos modelos permiten despliegue \textit{on-premise} sin dependencias externas.

En este proyecto se utiliza un modelo multilingüe de SentenceTransformers por compatibilidad con documentación en español/inglés y facilidad de despliegue local \cite{sbertDocs}.

% ============================================================
\section{Bases de Datos Vectoriales}
% ============================================================
La base vectorial define latencia y escalabilidad del sistema RAG. En grandes volúmenes, se usa búsqueda aproximada (ANN) con índices eficientes. Uno de los enfoques más utilizados es HNSW \cite{malkov2016hnsw}, un grafo navegable jerárquico que permite consultas rápidas a alta dimensionalidad.

\begin{table}[H]
\centering
\caption{Comparativa de bases de datos vectoriales evaluadas para el proyecto.}
\label{tab:vectordb-comparativa}
\begin{tabular}{lcccc}
\toprule
\textbf{Base de datos} & \textbf{Licencia} & \textbf{On-Premise} & \textbf{Filtrado} & \textbf{Híbrida} \\
\midrule
Qdrant   & Apache 2.0 & Sí & Sí & Sí \\
Weaviate & BSD-3 & Sí & Sí & Sí \\
Milvus   & Apache 2.0 & Sí & Sí & Sí \\
Chroma   & Apache 2.0 & Sí & Limitado & Limitado \\
Pinecone & Propietaria & No (SaaS) & Sí & Parcial \\
\bottomrule
\end{tabular}
\end{table}

Se selecciona \textbf{Qdrant} por ser \textit{on-premise}, por su filtrado por metadatos y por soporte en su ecosistema para búsqueda híbrida (dense + sparse) \cite{qdrantHybridTutorial,qdrantSparseVectors}. Además, dispone de API REST clara, lo que simplifica la integración con backend Python.

% ============================================================
\section{Cuantización y formatos para inferencia local}
\label{sec:cuantizacion}
% ============================================================
Para ejecutar LLMs localmente con recursos limitados se recurre a cuantización: reducir precisión de los pesos para disminuir VRAM y acelerar inferencia. En el ecosistema local, el formato GGUF se utiliza ampliamente para empaquetar modelos cuantizados junto con metadatos, especialmente en \texttt{llama.cpp} \cite{ggmlGgufSpec,hfGgufDocs}.

En este proyecto, la ejecución local se apoya en herramientas que consumen modelos GGUF y los sirven por API. Por ejemplo, Ollama documenta flujos de importación de modelos GGUF \cite{ollamaImportGguf}.

% ============================================================
\section{Ajuste fino eficiente con LoRA}
\label{sec:lora}
% ============================================================
El ajuste fino adapta un modelo generalista a un dominio concreto, pero entrenar todos los parámetros es costoso. LoRA propone entrenar adaptadores de bajo rango insertados en capas del modelo, reduciendo el número de parámetros actualizados y haciendo viable el entrenamiento con GPUs moderadas \cite{hu2021lora}:

\begin{equation}
W' = W + \Delta W = W + BA
\end{equation}

donde $A$ y $B$ son matrices de rango reducido y $r \ll \min(d,k)$.

% ============================================================
\section{Model Context Protocol (MCP)}
% ============================================================
MCP es un estándar abierto propuesto por Anthropic (2024) para integrar herramientas externas con agentes basados en LLM mediante una interfaz cliente-servidor tipada \cite{anthropicMcpAnnouncement}. Un servidor MCP expone herramientas (\textit{tools}) y un cliente MCP descubre e invoca dichas herramientas con argumentos estructurados.

En este proyecto se desarrolla un servidor MCP orientado a consulta normativa del BOE. Este servidor se ha validado en ejecución por terminal con un LLM y un cliente MCP compatible. Sin embargo, en el alcance actual, la integración MCP nativa no se incorpora en la interfaz web utilizada, por limitaciones del cliente. Por ello, la funcionalidad BOE dentro del sistema se integra mediante llamadas a API, manteniendo la operatividad y dejando preparada la migración a MCP completo en iteraciones futuras \cite{anthropicMcpAnnouncement}.

% ============================================================
\section{Tecnologías de soporte del ecosistema}
% ============================================================
Para implementar un RAG completo en local hacen falta piezas de integración:

\begin{itemize}
    \item \textbf{FastAPI:} backend en Python con documentación OpenAPI automática \cite{fastapiFeatures}.
    \item \textbf{LangChain / LlamaIndex:} utilidades para \textit{retrievers}, \textit{chunking} y orquestación RAG \cite{langchainRetrievalDocs,llamaindexRagDocs}.
    \item \textbf{LiteLLM:} proxy que unifica APIs de múltiples proveedores bajo un formato compatible con OpenAI \cite{litellmDocs}.
    \item \textbf{Ollama:} ejecución local de modelos y exposición por API, compatible con flujos GGUF \cite{ollamaImportGguf}.
\end{itemize}
